\section{Straightedge, compass and a nine-point circle}
\label{sec:constructions}
A line is given by $(A,u)$ with $N(u)>0$ and consists of $A+tu$.
This is an extensional representation: two such records describe the same
line precisely when their directions are proportional and one base point
lies on the other line. A circle $(C,s)$ uses its squared radius $s\ge0$,
with equation $N(X-C)=s$; $s=0$ gives a single point.
Every line can also be specified by two distinct points; every
circle can be specified by its center and a point on it.
All choices in the following classification are made on exact codes.

\paragraph{Two lines.}
For $A+tu=B+hv$ put $d=[u,v]$. If $d\ne0$, crossing with $v$
gives $t=[B-A,v]/d$ and crossing with $u$ gives $h=[B-A,u]/d$.
The two coordinates agree because their difference has zero cross
product with each vector of a basis. This proves existence and uniqueness.
If $d=0$, compare $[B-A,u]$ to zero to distinguish equality of lines
from disjoint parallel lines. The unique parallel to $(A,u)$ through $X$
is $(X,u)$, and its unique perpendicular there is $(X,Ju)$:
the solutions of $w\cdot u=0$ are the multiples of $Ju$.

\paragraph{A line and a circle.}
The foot from $C$ to $(A,u)$ is
\[
        F=A+u\frac{(C-A)\cdot u}{N(u)} .
\]
Indeed $F$ belongs to the line and $(C-F)\cdot u=0$; subtraction
gives uniqueness. For any point $F+tu$ on the line,
$N(F+tu-C)=N(F-C)+t^2N(u)$. Put $d=s-N(F-C)$.
There are zero, one or two intersections according as $d<0$, $d=0$
or $d>0$. They are $F$ at tangency and
$F\pm u\sqrt{d/N(u)}$ in the secant case. These exhaust the
solutions since $t^2=a^2$ implies $t=a$ or $t=-a$ in this field with decidable equality.
The formulas include the point-circle case.

\paragraph{Two circles.}
For $(C,s),(D,t)$ put $u=D-C$. If $u=0$, equality of squared radii
gives the same locus; different radii give disjoint loci. An identical
point-circle has just its center; an identical positive-radius circle
is returned as a whole circle. If $N(u)>0$, subtract the equations.
Every common point is on the line
\[
 C+k u+hJu,\quad k=\frac{N(u)+s-t}{2N(u)},\quad h\in K.
\]
Indeed this line is precisely the solution of
$2(X-C)\cdot u=N(u)+s-t$. Its substitution in the first circle gives
$h^2=s/N(u)-k^2$. The sign and square root give the exact number and
coordinates as in the preceding construction. Any candidate satisfying
the first circle and the subtraction satisfies the second circle.
All tests, including tangency, are primitive recursive.

\paragraph{Which of these witnesses can be drawn?}
Fix a frame with unit length and initial points. Let $E_0$ be the
subfield of $K$ generated by their coordinates in that frame.
Allow an intersection of already specified objects to return one of
its finitely many points, with either choice allowed. Parallel and
perpendicular constructions can be done with the two instruments.
For a perpendicular through $X$ to line $AB$, try the circles centered
at $X$ through $A$ or $B$. At least one gives a secant, since only the
foot on the line can give a tangent or a point-circle with its center on
the line. Its two intersection points have a perpendicular bisector
through $X$. To draw that bisector, use circles centered at the two
points through each other. Their two intersections lie on the bisector
and are distinct (their height squared is three quarters of the squared
chord length). Drawing two perpendiculars gives a parallel. Intersect
the bisector with the chord to obtain its midpoint. Equal input points have their common point as midpoint.

A point is obtainable by a finite construction exactly when both its
coordinates belong to the union of finite towers
\[
 E_0\subseteq E_1\subseteq\cdots\subseteq E_n,
 \qquad E_{i+1}=E_i(\sqrt{a_i}),\quad a_i\in E_i,\quad a_i>0.
\]
For necessity use the intersection formulas at each instruction.
For sufficiency, the constructions of Section~\ref{sec:frame-inverse}
realize addition and multiplication on the baseline. In detail, to copy
a displacement from $A$ to $D$ onto $C$, reflect $A$ through the midpoint
of $D,C$; if $A$ is that midpoint, return it. Otherwise the line
through $A$ and the midpoint and the circle about the midpoint through
$A$ give the other endpoint. The transfer diagram uses a parallel through
the input and its intersection with the other axis; the direction has
a nonzero transverse component. The coordinates of each initial point are obtained by parallels to the axes and this transfer in reverse. The product uses a parallel to the line
from $(1,0)$ to $(0,b)$, so it is also defined for $a=0$ or $b=0$.
Reflecting through the origin gives the negative. For $b\ne0$,
use the parallel through $(0,1)$ to the direction $(-1,b)$;
it meets the baseline at $(1/b,0)$. To construct $\sqrt a$, intersect
the vertical axis with the circle having diameter $(-1,0),(a,0)$.
At abscissa zero the circle equation is $y^2=a$; take the nonnegative
point and transfer its ordinate. Finally combine a horizontal and a
vertical displacement. A finite tower consists of field expressions
with finitely many such square roots since adjoining a root adjoins all
quotients of polynomial expressions in it. Two towers can be concatenated
with their radicands read in the enlarged intermediate fields.
This proves the stated characterization. It concerns generated finite
constructions with no arbitrary extra input points.

For example $T^3-2$ has a positive root in $K$. Over the unit rational
frame it cannot be drawn by these instruments. A rational root of the
monic integer cubic would have to be an integer divisor of $2$, and
none is a root, so its degree is three. Each tower step has degree one
or two, by reduction modulo $T^2-a_i$; when no square root is present
the quadratic is irreducible. Multiplying dimensions along finite field
extensions is justified by taking products of bases. Hence a degree-three
subfield cannot occur in such a tower, since three cannot divide a
power of two. The formula demanding a cube root is expressible with
points by the framed product graph. Its witness illustrates precisely
what the first-order completion adds to straightedge and compass.

\subsection{A theorem using incidence, length and construction}
\begin{proposition}[The nine-point circle, including right triangles]
For every noncollinear triangle the three side midpoints, the three
altitude feet and the three midpoints from its vertices to the
orthocenter lie on one explicitly constructible circle. Its radius is
half the circumradius.
\end{proposition}
\begin{proof}
The perpendicular bisectors of $AB$ and $AC$ have independent directions,
since $B-A,C-A$ do. Their unique intersection $O$ is equidistant from
all three vertices, and is their unique circumcenter. Translate $O$
to the origin; write the vertex vectors as $a,b,c$ with common norm
length $R>0$. Put $h=a+b+c$. Then
$(h-a)\cdot(c-b)=(b+c)\cdot(c-b)=0$; cycling proves that $h$ is on
each altitude. Two altitude directions are independent, so $h$ is their
unique common point. Put $n=h/2$. We will prove $|2p-h|=R$ for
each of the nine specified points $p$.

For $p=(b+c)/2$ one has $2p-h=-a$. For $p=(a+h)/2$ one has
$2p-h=a$. For the foot from $a$ to $bc$, let $m=(b+c)/2$ and
let $\Pi$ denote orthogonal projection on the direction $c-b$ through
zero. Because $m\cdot(c-b)=0$, the foot is $p=m+\Pi a$.
Thus $2p-h=2\Pi a-a$. The perpendicular summands $\Pi a$ and
$a-\Pi a$ have the same squared sum when the second is negated.
Consequently $|2\Pi a-a|=|a|=R$. Cyclic permutation proves the other
six assertions and gives $|p-n|=R/2$ throughout. The side midpoints
are noncollinear since their differences are half of two independent
side vectors, so the circle is unique. If the triangle is right several
of the nine points coincide; the argument only used nonzero side
lengths and a nonzero triangle determinant.
\end{proof}

The same determinant measures area: the oriented area of $ABC$ is
$[B-A,C-A]/2$. Splitting at any $X$ replaces this expression by the sum
for $XBC,AXC,ABX$ by bilinearity and alternating signs. Consequently a
triangulated polygon has signed area one half of the sum of
$[P_i,P_{i+1}]$ on its oriented boundary: the two opposite orientations
of each internal edge cancel. This definition uses only a finite sum.
